\section{Introduction}

Classifying an OWL~2~DL ontology means computing, for every pair of named
classes, whether the ontology entails that one is a subclass of the other,
together with the unsatisfiable classes and the consistency of the ontology.
The input is a finite set of axioms in the description logic $\mathcal{SROIQ}$
with datatypes: Boolean constructors, existential and universal restrictions,
role hierarchies and chains, inverse roles, nominals, and qualified number
restrictions.  The output is the entailed named-class taxonomy under the OWL~2
direct semantics \cite{w3c2012direct}.  Ontology editors and applications ask
further questions of the same kind: the types of each individual, the
individuals that are equal or different, the values of each property, and the
hierarchies of object and data properties.  The problem is decidable and
N2ExpTime-complete, and biomedical ontologies with hundreds of thousands of
axioms are large enough that a procedure needs extensive engineering to answer
within seconds.

Three families of procedures address this problem.  Tableau and hypertableau
reasoners such as FaCT++, Pellet, HermiT, and Konclude construct models and
refute candidate subsumptions; they handle the full language, and their
practical performance rests on absorption, blocking, caching, and dependency
tracking \cite{tsarkov2006factpp,sirin2007pellet,motik2009hypertableau,
steigmiller2014konclude}.  Consequence-based reasoners derive subsumptions by
saturating clauses.  ELK does this for the tractable \EL{} profile in
polynomial time \cite{kazakov2014elk}, and Sequoia extends the approach to
disjunctions, number restrictions, inverse roles, and nominals through a
context calculus \cite{bate2016sequoia,tenacucala2018consequence}.  Modular
and portfolio approaches such as MORe split an ontology between a tractable
and an expressive reasoner \cite{armasromero2012more}.  Each family has
ontologies on which it is the only practical choice, and the ORE competitions
showed that no single system dominates \cite{parsia2017ore}.

A user who wants the best procedure for each ontology therefore has to install
several reasoners and choose among them.  The reasoners differ in input
syntaxes, in the queries they answer, and in how they report an input they
cannot handle; ELK, for example, classifies an ontology outside \EL{} by
ignoring the axioms it cannot represent.  Several of the systems are barely
maintained \cite{abicht2023owl}, and most run on the Java virtual machine,
whose start-up time and memory dominate the cost of classifying a small
ontology.  Correctness is established by agreement between systems, and the
systems disagree on a measurable fraction of real ontologies
\cite{lee2015correctness,parsia2017ore}.

\KM{} (KM) addresses these limitations in one system.  KM is a single Rust
executable that contains an \EL{} completion procedure, a disjunctive
consequence-based context calculus with nominals, a native hypertableau
procedure, and a Rust port of Konclude's completion and classification
algorithms.  A router selects a procedure from features of the source
ontology.  Every procedure follows one contract: it returns a complete answer
or it declines, and a decline passes the input to a fallback procedure, so an
input outside the fragment of a fast procedure is answered by a slower one
instead of being approximated.  KM refuses an input that violates the OWL~2~DL
global restrictions and reports the violation.  For the correctness of its
answers KM adds a second line of evidence to agreement with other reasoners:
a Lean~4 development \cite{demoura2021lean4} states the OWL~2 direct
semantics, proves the reasoning rules sound, and provides checkers that accept
or reject the evidence a KM run emits for its answer.

Beyond classification, KM computes the types and equalities of individuals,
property values, the object and data property hierarchies, justifications for
entailments, a finite model of a consistent ontology, and incremental
reclassification after an edit.  A Prot\'eg\'e plugin, which bundles the
executable for five platforms, answers the calls of the OWL API reasoner
interface that HermiT answers.  Almost all of KM's code and proofs were
written by large-language-model coding agents under human direction; this
paper describes the resulting system and its evaluation.

We evaluate KM v1.7.0 on the ORE~2015 corpus \cite{parsia2017ore} under a
protocol of one core, 240 seconds, and 20~GiB per ontology, with correctness
decided by agreement with Konclude, HermiT, Openllet, and JFact.  On
\HOInputs{} held-out ontologies KM solves \HOKMSolved{} and Konclude
\HOKCSolved{}, and KM is faster and uses less memory in the mean and in the
median on the inputs both solve.  We compare KM with OxidOWL
\cite{sieve2026oxidowlcodingagents}, a hypertableau reasoner for OWL~2~DL
that is also written in Rust, under the same protocol; OxidOWL finishes
classification on \OXFinished{} of these ontologies and returns the reference
hierarchy on \OXHierarchyCorrect.  We further report
differential tests against HermiT on generated ontologies, a call-by-call
comparison of the Prot\'eg\'e plugin with HermiT, and the fraction of answers
that the Lean checkers certify.

KM is available at
\url{https://github.com/bio-ontology-research-group/kobayashi-marust} under
the BSD 3-Clause licence, with the Konclude port under the LGPL~v3.
Section~\ref{sec:background} reviews the procedures KM combines.
Sections~\ref{sec:overview} to~\ref{sec:routing} describe the architecture,
the procedures, and the router; Section~\ref{sec:services} describes the
reasoning services and interfaces; Section~\ref{sec:certification} describes
the Lean development; and Section~\ref{sec:evaluation} reports the
evaluation.
